\section{Background}
\label{sec:background}

\paragraph{Diffusion Models}\citep{SohlDickstein2015DeepUL,Song2020ScoreBasedGM,ho2020denoising} generate data by approximating the reverse ODE to a stochastic forward process which transforms data to noise. They have become the standard approach for generative modeling of  images~\cite{dhariwal2021diffusion,ramesh2022hierarchical,saharia2022photorealistic,Rombach_2022,balaji2022ediffi} and videos~\citep{singer2022makeavideo,ho2022imagen,esser2023structure, blattmann2023align,gupta2023photorealistic}. Since these models can be derived both via a variational lower bound on the negative likelihood~\citep{SohlDickstein2015DeepUL} and score matching~\citep{Hyvrinen2005EstimationON,Vincent2011ACB,song2020generative}, 
%
various formulations of forward- and reverse processes~\citep{Song2020ScoreBasedGM,dockhorn2021score}, model parameterizations~\citep{ho2020denoising,ho2022classifierfree,Karras2022ElucidatingTD}, loss weightings~\citep{ho2020denoising, Karras2022ElucidatingTD} and ODE 
solvers~\citep{song2022denoising,lu2023dpmsolver, dockhorn2022genie} have led to a large number of different training objectives and sampling procedures. More recently, the seminal works of \citet{Kingma2023UnderstandingDO} and \citet{Karras2022ElucidatingTD} have proposed unified formulations and introduced new theoretical and practical insights for training~\cite{Karras2022ElucidatingTD,Kingma2023UnderstandingDO} and inference~\cite{Karras2022ElucidatingTD}. However, despite these improvements, the trajectories of common ODEs involve partly significant amounts of curvature~\citep{Karras2022ElucidatingTD,liu2022flow}, which requires increased amounts of solver steps and, thus, renders fast inference difficult. To overcome this, we adopt rectified flow models whose formulation allows for learning straight ODE trajectories.  

\paragraph{Rectified Flow Models}
\cite{liu2022flow,albergo2022building,lipman2023flow} approach generative modeling by
constructing a transport map between two distributions through an ordinary
differential equation (ODE). This approach has close connections to continuous
normalizing flows (CNF) \cite{Chen2018NeuralOD} as well as diffusion models.
Compared to CNFs, Rectified Flows and Stochastic Interpolants have the
advantage that they do not require simulation of the ODE during training.
Compared to diffusion models, they can result in ODEs that are faster to
simulate than the probability flow ODE \cite{Song2020ScoreBasedGM} associated
with diffusion models. Nevertheless, they do not result in optimal transport
solutions, and multiple works aim to minimize the trajectory curvature further
\cite{lee2023minimizing,alex2023improving,aramalex2023multisample}.
\cite{dao2023flow,ma2024sit} demonstrate the feasibility of rectified flow
formulations for class-conditional image synthesis, \cite{fischer2023boosting} for
latent-space upsampling, and \cite{liu2023instaflow}
apply the reflow procedure of \cite{liu2022flow} to distill a pretrained
text-to-image model \cite{Rombach_2022}. Here, we are interested in rectified
flows as the foundation for text-to-image synthesis with fewer sampling steps.
We perform an extensive comparison between different formulations and loss
weightings and propose a new timestep schedule for training of
rectified flows with improved performance.



\paragraph{Scaling Diffusion Models}
The transformer architecture~ \cite{vaswani2017attention} is well known for its scaling properties in NLP~\cite{kaplan2020scaling} and computer vision tasks~\cite{dosovitskiy2020image,zhai2022scaling}. For diffusion models, U-Net architectures~\cite{Ronneberger_2015} have been the  dominant choice~\cite{ho2020denoising,Rombach_2022,balaji2022ediffi}.
While some recent works explore diffusion transformer backbones~\cite{Peebles_2023,chen2023pixart,ma2024sit}, scaling laws for text-to-image diffusion models remain unexplored.

%

%
%
%
%
%
%
%
%
%

%
%
%
%
%
%

